% M1596
\section*{Заяви та декларації}
% M1597
\subsection*{Внески авторів}
Тетяна Поліщук аналітично розв’язала задачу під керівництвом Михайла Дудика. Михайло Дудик брав участь у розробленні концепції та методики й керував аналітичною роботою. Анатолій Камінський рецензував науковий зміст рукопису. Михайло Селіванов незалежно перевірив і виправив аналітичні й числові результати, відтворив і перевірив розрахунки, підготував рисунки, брав участь в інтерпретації результатів, відновив і відредагував рукопис, підготував англомовну версію та координував подання. Усі автори переглянули остаточний рукопис.

% M1608
\subsection*{Фінансування}
Дослідження не отримувало зовнішнього фінансування. Роботу виконано в межах базового фінансування установ авторів.

% M1612
\subsection*{Доступність даних}
Усі дані, які підтверджують результати дослідження, наведено в статті. Числові криві й табличні значення отримано за аналітичними співвідношеннями, наведеними в рукописі.

% M1617
\subsection*{Конфлікт інтересів}
Автори заявляють про відсутність конфлікту інтересів.

% M1620
\appendix

enewcommand{\theHequation}{appendix.\arabic{equation}}
% M1621
\section{Поля напружень і переміщень поблизу кутової точки}
% M1622
\label{app:asymptotics}
% M1623
\providecommand{\SingularityExponentFigure}{рис.~\ref{fig:singularity-exponents}}

% M1625
\setcounter{equation}{0}
% M1626
\renewcommand{\theequation}{A.\arabic{equation}}

% M1628
У цьому додатку зібрано характеристичні співвідношення та коефіцієнти асимптотичних полів, використані в основному тексті. Для плоскої деформації безрозмірне відношення пружних характеристик $e$ і параметри матеріалів $\kappa_i$ мають вигляд
% M1631
\begin{equation*}
  e=\frac{1+\nu_2}{1+\nu_1}\frac{E_1}{E_2},
  \qquad
  \kappa_i=3-4\nu_i,\qquad i=1,2.
\end{equation*}
% M1636
Для кожного характеристичного показника комплексна змінна перетворення Мелліна $p$ та узагальнене позначення показника сингулярності $\lambda_*$ пов’язані співвідношенням
% M1638
\begin{equation*}
  p=-1-\lambda_*,
  \qquad -1<\lambda_*<0.
\end{equation*}
% M1642
За фізичний показник беремо найменший дійсний корінь в інтервалі $(-1,0)$, неперервно пов’язаний із відповідною гілкою розв’язку. Випадок $\alpha=\pi/2$ відповідає виродженій границі прямолінійної межі поділу й не розглядається як звичайна внутрішня точка характеристичної гілки.

% M1648
\subsection{Кутова точка непошкодженої межі поділу}

% M1650
Для конфігурації без міжфазної тріщини зсуву сингулярна частина поля напружень на межі поділу має вигляд~\citep{NazarenkoKipnis2022Stress}
% M1652
\begin{equation}
\label{eq:app-intact-field}
\begin{aligned}
  \tau_{r\theta}(r,\alpha)
    &=Cg_1r^{\lambda_0}+o\!\left(r^{\lambda_0}\right),\\
  \sigma_\theta(r,\alpha)
    &=Cg_2r^{\lambda_0}+o\!\left(r^{\lambda_0}\right),
    \qquad r\to0,
\end{aligned}
\end{equation}
% M1662
де $C$ — амплітуда локального сингулярного поля, $\lambda_0$ — показник сингулярності в кутовій точці непошкодженої межі поділу, а $g_1$ і $g_2$ — безрозмірні кутові коефіцієнти, що дорівнюють
% M1665
\begin{equation*}
\begin{aligned}
  g_1={}&(1-e)\sin[\lambda_0(\pi-\alpha)]
  \bigl[(\lambda_0+1)\sin2\alpha
       +\sin\{2(\lambda_0+1)\alpha\}\bigr]\\
  &+e(1+\kappa_2)\sin(\lambda_0\pi)
      \sin[(\lambda_0+2)\alpha],\\[2pt]
  g_2={}&(1-e)\cos[\lambda_0(\pi-\alpha)]
  \bigl[(\lambda_0+1)\sin2\alpha
       +\sin\{2(\lambda_0+1)\alpha\}\bigr]\\
  &+e(1+\kappa_2)\sin(\lambda_0\pi)
      \cos[(\lambda_0+2)\alpha].
\end{aligned}
\end{equation*}
% M1679
Показник $\lambda_0$ визначається характеристичним рівнянням
% M1680
\begin{equation}
\label{eq:app-D0-root}
  D_0(-1-\lambda_0)=0,
  \qquad -1<\lambda_0<0,
\end{equation}
% M1685
де визначник $D_0(p)$ для кутової точки непошкодженої межі має вигляд
% M1686
\begin{equation*}
  D_0(p)=b_0(p)+e\,b_1(p)+e^2b_2(p),
\end{equation*}
% M1689
а допоміжні функції $b_0(p)$, $b_1(p)$ й $b_2(p)$ задано формулами
% M1690
\begin{equation*}
\begin{aligned}
  b_0(p)={}&
  [\sin(2p\alpha)+p\sin(2\alpha)]
  [\kappa_1\sin\{2p(\pi-\alpha)\}+p\sin(2\alpha)],\\[2pt]
  b_1(p)={}&(1+\kappa_1)(1+\kappa_2)\sin^2(p\pi)\\
  &-[\sin(2p\alpha)+p\sin(2\alpha)]
    [\kappa_1\sin\{2p(\pi-\alpha)\}+p\sin(2\alpha)]\\
  &-[\sin\{2p(\pi-\alpha)\}-p\sin(2\alpha)]
    [\kappa_2\sin(2p\alpha)-p\sin(2\alpha)],\\[2pt]
  b_2(p)={}&
  [\sin\{2p(\pi-\alpha)\}-p\sin(2\alpha)]
  [\kappa_2\sin(2p\alpha)-p\sin(2\alpha)].
\end{aligned}
\end{equation*}
% M1705
\subsection{Кутова точка зі стаціонарною V-подібною міжфазною тріщиною зсуву}

% M1707
Для заданої стаціонарної V-подібної міжфазної тріщини зсуву кожна з двох симетричних гілок має довжину $l$. На відстанях $r\ll l$ нормальне напруження вздовж бісектриси сектора, зайнятого матеріалом~$i$, має асимптотичний вигляд~\citep{NazarenkoKipnis2022Stress}
% M1711
\begin{equation}
\label{eq:app-cracked-field}
  \sigma_\theta(r,\beta_i)
  =CQ_i r^\lambda+o\!\left(r^\lambda\right),
  \qquad r\to0,
  \qquad
  \beta_1=\pi,\quad \beta_2=0.
\end{equation}
% M1719
Тут $\lambda$ — показник сингулярності в кутовій точці конфігурації з тріщиною, а $\beta_i$ — напрямок відповідної бісектриси. Амплітудний коефіцієнт $Q_i$ має вигляд
% M1722
\begin{equation*}
  Q_i=g_1q_i l^{\lambda_0-\lambda},
\end{equation*}
% M1725
де безрозмірний коефіцієнт $q_i$ і похідна $D'(p)$ визначника для кутової точки з тріщиною дорівнюють
% M1727
\begin{equation*}
  q_i=
  \frac{2S_i(-1-\lambda)K^+(-1-\lambda)G^+(-1-\lambda)}
  {(\lambda-\lambda_0)D'(-1-\lambda)
   K^+(-1-\lambda_0)G^+(-1-\lambda_0)},
  \qquad
  D'(p)=\frac{\dd D(p)}{\dd p}.
\end{equation*}
% M1735
Функції $S_i(p)$ — допоміжні функції чисельника, визначені нижче. Множники $K^+(p)$ і $G^+(p)$ задано далі в цьому підрозділі. Плюс-фактор базового ядра $G^+(p)$ слід відрізняти від $G_i^+(p)$, який бере участь у факторизації ядра $G_i(p)$ для зони передруйнування в розділі~\ref{sec:derivation}. У явному вигляді маємо
% M1740
\begin{equation*}
\begin{aligned}
  S_1(p)={}&2e(1+\kappa_2)s_{11}(p)\sin[(p+1)\alpha]\\
  &+s_{12}(p)
  \Bigl[(1+\kappa_1)\sin\{p(\pi-\alpha)-\alpha\}
        +2(1-e)s_{13}(p)\Bigr],\\[2pt]
  S_2(p)={}&s_{21}(p)
  \Bigl[e(1+\kappa_2)\sin[(p+1)\alpha]
        +2(1-e)s_{22}(p)\Bigr]\\
  &-2(1+\kappa_1)\sin\{p(\pi-\alpha)-\alpha\}\,s_{23}(p),
\end{aligned}
\end{equation*}
% M1752
де допоміжні функції $s_{mn}(p)$, при $m=1,2$ і $n=1,2,3$, мають вигляд
% M1754
\begin{equation*}
\begin{aligned}
  s_{11}(p)&=p\cos(p\pi-\alpha)\sin\alpha
             -\sin[p(\pi-\alpha)]\cos(p\alpha),\\
  s_{12}(p)&=p\sin2\alpha+\sin(2p\alpha),\\
  s_{13}(p)&=p\cos[p(\pi-\alpha)]\sin\alpha
             -\sin[p(\pi-\alpha)]\cos\alpha,\\
  s_{21}(p)&=\sin[2p(\pi-\alpha)]-p\sin2\alpha,\\
  s_{22}(p)&=p\cos(p\alpha)\sin\alpha
             +\sin(p\alpha)\cos\alpha,\\
  s_{23}(p)&=p\cos(p\pi+\alpha)\sin\alpha
             +\cos[p(\pi-\alpha)]\sin(p\alpha).
\end{aligned}
\end{equation*}
% M1768
Показник сингулярності $\lambda$ у кутовій точці конфігурації із заданою стаціонарною V-подібною міжфазною тріщиною задовольняє рівняння
% M1770
\begin{equation}
\label{eq:app-D-root}
  D(-1-\lambda)=0,
  \qquad -1<\lambda<0,
\end{equation}
% M1775
де характеристичний визначник $D(p)$ має вигляд
% M1776
\begin{equation*}
  D(p)=a_0(p)+e\,a_1(p),
\end{equation*}
% M1779
а допоміжні функції $a_0(p)$ й $a_1(p)$ задано формулами
% M1780
\begin{equation*}
\begin{aligned}
  a_0(p)={}&(1+\kappa_1)
  [\cos\{2p(\pi-\alpha)\}-\cos2\alpha]
  [\sin(2p\alpha)+p\sin2\alpha],\\
  a_1(p)={}&(1+\kappa_2)
  [\cos(2p\alpha)-\cos2\alpha]
  [\sin\{2p(\pi-\alpha)\}-p\sin2\alpha].
\end{aligned}
\end{equation*}
% M1790
У виразах для коефіцієнтів $q_i$ використано базове ядро
% M1791
\begin{equation*}
  G(p)=
  \frac{(e+\kappa_1)(1+e\kappa_2)\sin(\pi p)D(p)}
  {(e+\kappa_1+1+e\kappa_2)\cos(\pi p)D_0(p)}
\end{equation*}
% M1796
і його плюс-фактор
% M1797
\begin{equation*}
  G^+(p)=
  \exp\!\left\{
    \frac{1}{2\pi\ii}
    \int_{-\ii\infty}^{\ii\infty}
    \frac{\log G(z)}{z-p}\,\dd z
  \right\},
  \qquad \operatorname{Re}p<0.
\end{equation*}
% M1806
Тут $z$ — комплексна змінна інтегрування, $\ii^2=-1$, а контур спрямовано вгору вздовж уявної осі. Гамма-функціональний множник має вигляд
% M1808
\begin{equation*}
  K^+(p)=\frac{\Gamma(1-p)}{\Gamma(\tfrac12-p)}.
\end{equation*}
% M1811
Тут $\Gamma$ — гамма-функція Ейлера. Для дійсного $p<0$ величину $G^+(p)$ рівнозначно можна обчислити за допомогою дійсного інтеграла
% M1813
\begin{equation*}
  \log G^+(p)
  =-\frac{p}{\pi}\int_0^\infty
    \frac{\log G(\ii t)}{t^2+p^2}\,\dd t.
\end{equation*}
% M1818
Тут $t$ — дійсна змінна інтегрування.

% M1820
\subsection{Кутова точка із зоною передруйнування}

% M1822
Для зони передруйнування в матеріалі~$i$ ядро функціонального рівняння з розділу~\ref{sec:derivation} має вигляд
% M1824
\begin{equation*}
  G_i(p)=\frac{D_i(p)\cos(\pi p)}{D(p)\sin(\pi p)}.
\end{equation*}
% M1827
Відповідні характеристичні визначники дорівнюють
% M1828
\begin{equation}
\label{eq:app-post-zone-determinants}
\begin{aligned}
  D_1(p)={}&2e(1+\kappa_2)
  [\cos(2p\alpha)-\cos2\alpha]
  \{\sin^2[p(\pi-\alpha)]-p^2\sin^2\alpha\}\\
  &+(1+\kappa_1)
  [\sin(2p\alpha)+p\sin2\alpha]
  [\sin\{2p(\pi-\alpha)\}-p\sin2\alpha],\\[3pt]
  D_2(p)={}&2(1+\kappa_1)
  [\cos\{2p(\pi-\alpha)\}-\cos2\alpha]
  [\sin^2(p\alpha)-p^2\sin^2\alpha]\\
  &+e(1+\kappa_2)
  [\sin\{2p(\pi-\alpha)\}-p\sin2\alpha]
  [\sin(2p\alpha)+p\sin2\alpha].
\end{aligned}
\end{equation}
% M1845
Показник $\lambda_i$ для конфігурації із зоною в матеріалі~$i$ визначається рівнянням
% M1847
\begin{equation}
\label{eq:app-post-zone-root}
  D_i(-1-\lambda_i)=0,
  \qquad -1<\lambda_i<0,
  \qquad i=1,2.
\end{equation}
% M1853
Ці корені визначають криві $\lambda_1$ і $\lambda_2$ на \SingularityExponentFigure. Вони характеризують сингулярність кутової точки за наявності зони передруйнування. У формулах для $d_i$, $\delta_i$ й $J_i$ використовують $\lambda$ з рівняння~\eqref{eq:app-D-root}.

% M1860
\begin{english}\bibliography{references}\end{english}
